\documentclass[12pt]{article}
\usepackage[utf8]{inputenc}
\usepackage[T1]{fontenc}
\usepackage[margin=1in]{geometry}
\usepackage{setspace}
\usepackage{xcolor}
\usepackage[colorlinks=true, urlcolor=blue!60!black, linkcolor=blue!60!black]{hyperref}
\setlength{\parindent}{0em}
\setlength{\parskip}{0.6em}
\renewcommand{\baselinestretch}{1.0}
\begin{document}
\begin{center}
{\LARGE \textit{Chae Chan Ping v. United States}, 130 U.S. 581 (1889)\par}
\vspace{0.5em}
{\large Case Brief --- Two-Page Edition\par}
\end{center}

\section{Memory Jogger}
\textit{The Chinese Exclusion Case}: a later statute overrides a conflicting treaty, and Congress's inherent sovereign power to exclude noncitizens is conclusive on the courts (\href{https://supreme.justia.com/cases/federal/us/130/581/}{130 U.S. 581}).

\section{Detailed Case Facts}
Chae Chan Ping, a Chinese laborer living lawfully in San Francisco, left for China in 1887 with a reentry certificate issued under the Chinese Exclusion Act, as treaties guaranteed resident laborers the right to ``go and come of their own free will and accord.'' While he was at sea returning, Congress passed the Scott Act of 1888, voiding all certificates, and he was refused permission to land.

\section{Procedural History}
The federal circuit court denied his habeas petition; the Supreme Court affirmed.

\section{Judicial Votes}
Affirmed, 8--0. Justice Field wrote for a unanimous Court, with one seat vacant.

\section{Holding}
Statutes and treaties are equally supreme, so a later statute abrogates a conflicting treaty as domestic law. The power to exclude noncitizens is an inherent incident of sovereignty, and Congress's exclusion judgment is conclusive upon the Judiciary.

\section{Analysis of Opinions}
Field conceded the treaty breach but held that ``the last expression of the sovereign will must control,'' leaving diplomatic grievances to the political branches. He grounded the exclusion power not in an enumerated clause but in nationhood itself, invoking \textit{The Schooner Exchange} and \textit{Cohens v. Virginia}. In racialized terms, he held that if Congress deems ``foreigners of a different race'' who ``will not assimilate with us'' dangerous, its judgment is ``conclusive upon the judiciary.'' A reentry certificate was a license ``revocable at any time, at its pleasure.'' The plenary-power doctrine grew through \textit{Fong Yue Ting}, \textit{Knauff}, and \textit{Mezei}, was limited by \textit{Wong Wing} and \textit{Mandel}, and was applied in \textit{Trump v. Hawaii}.

\section{Examples: Future Applications}
\textbf{Same-Side:} a statute voiding visas promised by a visa-waiver treaty controls under the later-in-time rule. \textbf{Opposite-Side:} imprisoning noncitizens at hard labor without a jury violates the Fifth and Sixth Amendments (\textit{Wong Wing}). \textbf{Fence-Sitter:} a statute excluding immigrants of a named religion, given \textit{Trump v. Hawaii}'s emphasis on facial neutrality.

\section{Critique}
Critics say inherent sovereignty departs from \textit{McCulloch}'s enumerated powers and creates a constitutional black hole. Field's racial reasoning shows the danger of unchecked deference, which critics say \textit{Trump v. Hawaii} perpetuates.

\section{Key Quotations}
\begin{quote}
``For local interests the several states of the Union exist, but for national purposes, embracing our relations with foreign nations, we are but one people, one nation, one power. ... If, therefore, the government of the United States, through its legislative department, considers the presence of foreigners of a different race in this country, who will not assimilate with us, to be dangerous to its peace and security, their exclusion is not to be stayed because at the time there are no actual hostilities with the nation of which the foreigners are subjects. ... In both cases its determination is conclusive upon the judiciary.''
\end{quote}
\noindent --- Justice Field, Opinion of the Court, 130 U.S.\ at 606.

\end{document}
